\steppingstone{Can Every Multiplication Behave Like Growth?}
\label{stone:multiplying}

Growing by \(\varphi\) turned out to be wonderfully simple. Every short piece became one long piece, and every long piece became one short piece and one long piece. The geometry performed the work. PhiBase merely counted the result.

That success invites a new question.

Was growing special because one of the quantities happened to be \(\varphi\), or can every multiplication be understood in the same spirit?

Suppose we wish to multiply

\[
P(a,b)
\]

by

\[
P(c,d).
\]

Instead of trying to understand the whole product at once, let us pretend we are still growing. During growth, one quantity had only a single nonzero component. Can we arrange for the same thing to happen here?

That is the real problem.

We are looking for a companion that allows one part of the multiplication to disappear, leaving something as simple as growth.

For the moment, however, let us simply build the product.

Write the two quantities as

\[
P(a,b)=a+b\varphi,
\]

and

\[
P(c,d)=c+d\varphi.
\]

Now expand them exactly as you would any two binomials,

\[
(a+b\varphi)(c+d\varphi)
=
ac+ad\varphi+bc\varphi+bd\varphi^{2}.
\]

Only one unfamiliar term appears,

\[
\varphi^{2}.
\]

The previous stepping stone has already shown us what to do.

Replace it with

\[
\varphi+1.
\]

Now collect the short pieces together, then collect the long pieces.

\begin{center}
\begin{tabular}{rcl}
Short part &=& \(ac+bd\)\\[0.4ex]
Long part  &=& \(ad+bc+bd\)
\end{tabular}
\end{center}

The product is therefore

\[
P(a,b)P(c,d)
=
P(ac+bd,\;ad+bc+bd).
\]

The answer is another PhiBase quantity.

Nothing has escaped from the arithmetic.

Nothing has been approximated.

The geometry quietly returned every occurrence of \(\varphi^{2}\) to the same two building pieces with which we began.

\section*{One Example}

Multiply

\[
P(2,1)
\]

by

\[
P(3,2).
\]

The calculation is simply a matter of counting.

\begin{center}
\begin{tabular}{rcl}
Short part &=& \(2\cdot3+1\cdot2=8\)\\[0.6ex]
Long part  &=& \(2\cdot2+1\cdot3+1\cdot2=9\)
\end{tabular}
\end{center}

Therefore

\[
P(2,1)P(3,2)=P(8,9).
\]

The same result appears if the two quantities are expanded in ordinary algebra. PhiBase has not changed the mathematics. It has preserved the construction.

\section*{Builder's Notebook}

Multiply

\[
P(1,1)P(1,1),
\]

\[
P(2,3)P(1,2),
\]

and

\[
P(4,1)P(2,2).
\]

After each multiplication, identify the short part and the long part separately before writing the final PhiBase quantity.

\section*{Looking Ahead}

Growth became simple because one quantity contained only a single nonzero part.

Can every multiplication be made to look that simple?

If we could find the right companion for a PhiBase quantity, one whole component of the product might disappear.

That question leads directly to division.
